% ============ 第4章 §4.2 Hardy-Littlewood-Sobolev不等式(原书 pp.319-324) ============
\section{Hardy--Littlewood--Sobolev 不等式}\label{sec:4.2}
\index{Hardy--Littlewood--Sobolev 不等式}
\index{分数阶积分算子}


先复习 Young 不等式.

\noindent\textbf{Young 不等式}\quad 首先准备一个引理.

\begin{lemma}\label{lem:4.2.1}
设 $f(x)$ 为 $\mathbb{R}^{n}$ 上的可测函数.$f \in L^{p}(\mathbb{R}^{n})$($1 \leq p \leq \infty$)的充分必要条件为:设 $\frac{1}{p} + \frac{1}{p'} = 1$,存在满足
\begin{equation*}
\left|\int_{\mathbb{R}^{n}}f(x)g(x)\dd x\right| \leq M\|g\|_{L^{p'}(\mathbb{R}^{n})} \quad (g \in L^{p'}(\mathbb{R}^{n}))
\end{equation*}
的常数 $M$.特别地,此时
\begin{equation*}
\|f\|_{L^{p}(\mathbb{R}^{n})} \leq M
\end{equation*}
成立.\footnote{以后,对 $1 \leq p \leq \infty$,用 $p'$ 表示满足 $1/p + 1/p' = 1$ 的数,即 $p' = p/(p-1)$.}
\end{lemma}

\begin{proof}
只要设 $f, g$ 均为非负来证明即可.必要性由 Hölder 不等式显然,故证明充分性.

若 $p = 1$,令 $g \equiv 1 \in L^{\infty}$,则 $\|f\|_{1} \leq M$.

设 $1 < p < \infty$.令
\begin{equation*}
f_{t}(x) := f(x)\chi_{B(0,t)}(x)\chi_{[0,t]}(f(x)) \quad (t > 0),
\end{equation*}
则
\begin{equation*}
f_{t}(x)^{p-1} \in L^{p'}(\mathbb{R}^{n}) \quad (t > 0).
\end{equation*}
利用这一点,并注意 $p' = p/(p-1)$,
\begin{equation*}
\int_{\mathbb{R}^{n}}f_{t}(x)^{p}\dd x = \int_{\mathbb{R}^{n}}f_{t}(x)f_{t}(x)^{p-1}\dd x \leq M\left(\int_{\mathbb{R}^{n}}f_{t}(x)^{p}\dd x\right)^{1/p'}.
\end{equation*}
由于 $1/p + 1/p' = 1$,将两边除以
\begin{equation*}
\left(\int_{\mathbb{R}^{n}}f_{t}(x)^{p}\dd x\right)^{1/p'},
\end{equation*}
得
\begin{equation*}
\|f_{t}\|_{p} \leq M.
\end{equation*}
令 $t \to \infty$,由单调收敛定理得 $\|f\|_{p} \leq M$.

$p = \infty$ 时,令
\begin{equation*}
g_{y,t}(x) := \frac{\chi_{B(y,t)}(x)}{|B(y,t)|},
\end{equation*}
则
\begin{equation*}
\|g_{y,t}\|_{1} = 1,
\end{equation*}
故
\begin{equation*}
\int_{\mathbb{R}^{n}}f(x)g_{y,t}(x)\dd x \leq M.
\end{equation*}
令 $t \to 0$,左端由 Lebesgue 微分定理在几乎处处 $y \in \mathbb{R}^{n}$ 收敛到 $f(y)$.因此
\begin{equation*}
f(y) \leq M \quad (\mathrm{a.e.} \ y \in \mathbb{R}^{n}),
\end{equation*}
得 $\|f\|_{\infty} \leq M$.
\end{proof}

进入正题.

\begin{proposition}[Young 不等式]\label{prop:4.2.2}
设 $p, q, r$ 满足
\begin{equation}
\frac{1}{r} = 1 - \left(\frac{1}{p} - \frac{1}{q}\right), \quad 1 \leq p \leq q \leq \infty, \quad 1 \leq r \leq \infty.\tag{4.2.1}
\label{eq:4.2.1}\end{equation}
此时,令
\begin{equation*}
Tf(x) := \int_{\mathbb{R}^{n}}K(x, y)f(y)\dd y,
\end{equation*}
若假设
\begin{equation*}
\sup_{x \in \mathbb{R}^{n}}\left(\int|K(x, y)|^{r}\dd y\right)^{1/r} \leq C, \quad \sup_{y \in \mathbb{R}^{n}}\left(\int|K(x, y)|^{r}\dd x\right)^{1/r} \leq C,
\end{equation*}
则
\begin{equation}
\|Tf\|_{L^{q}(\mathbb{R}^{n})} \leq C\|f\|_{L^{p}(\mathbb{R}^{n})}\tag{4.2.2}
\label{eq:4.2.2}\end{equation}
成立.
\end{proposition}

\begin{proof}
用 M. Riesz 插值定理证明.即插值
\begin{equation*}
\|Tf\|_{L^{\infty}(\mathbb{R}^{n})} \leq C\|f\|_{L^{r'}(\mathbb{R}^{n})}
\end{equation*}
与
\begin{equation*}
\|Tf\|_{L^{r}(\mathbb{R}^{n})} \leq C\|f\|_{L^{1}(\mathbb{R}^{n})},
\end{equation*}
得到结论.

前一不等式由 Hölder 不等式可知成立.

后一不等式:设 $f \in L^{1}(\mathbb{R}^{n})$,$g \in L^{r'}(\mathbb{R}^{n})$,
\begin{equation}
\begin{aligned}
\left|\int Tf(x)g(x)\dd x\right| & = \left|\int\int K(x, y)f(y)g(x)\dd y\dd x\right| \\
& = \left|\int f(y)\int K(x, y)g(x)\dd x\dd y\right| \leq C\|f\|_{1}\times\|g\|_{r'},
\end{aligned}
\end{equation}
可以算出,由引理~\ref{lem:4.2.1} 可知成立.
\end{proof}

\noindent\textbf{分数阶积分算子}\quad 用 Hardy--Littlewood 极大函数的 $L^{p}(\mathbb{R}^{n})$ 估计来证明分数阶积分算子的 $L^{p}(\mathbb{R}^{n}) \to L^{q}(\mathbb{R}^{n})$ 有界性.

设 $r > 1$,分数阶积分算子 $I_{r}$ 定义为
\begin{equation*}
I_{r}f(x) := \int_{\mathbb{R}^{n}}\frac{f(y)}{|x-y|^{n/r}}\dd y \quad (r > 1).
\end{equation*}
关于分数阶积分算子,下面的定理成立.

\begin{theorem}[Hardy--Littlewood--Sobolev 不等式]\label{thm:4.2.3}
设 $p, q, r$ 满足
\begin{equation}
\frac{1}{r} = 1 - \left(\frac{1}{p} - \frac{1}{q}\right), \quad 1 < p < q < \infty, \quad 1 < r < \infty.\tag{4.2.3}
\label{eq:4.2.3}\end{equation}
此时,若 $f \in L^{p}(\mathbb{R}^{n})$,则
\begin{equation}
\|I_{r}f\|_{L^{q}(\mathbb{R}^{n})} \leq C_{n,p,q}\|f\|_{L^{p}(\mathbb{R}^{n})}\tag{4.2.4}
\label{eq:4.2.4}\end{equation}
成立.
\end{theorem}

Hardy--Littlewood--Sobolev 不等式与 Young 不等式的区别在于,$I_{r}$ 的核 $|y|^{-n/r}$ 不属于 $L^{r}(\mathbb{R}^{n})$.

\begin{proof}[定理~\ref{thm:4.2.3} 的证明]
固定 $f \in L^{p}(\mathbb{R}^{n})$ 与 $x \in \mathbb{R}^{n}$,估计 $I_{r}f(x)$.

定义
\begin{equation*}
f_{t}(x) := \frac{1}{t^{n}}\int_{B(x,t)}|f(y)|\dd y \quad (t > 0).
\end{equation*}
此时,对 $y \in \mathbb{R}^{n}$($y \neq x$),
\begin{equation*}
\int_{0}^{\infty}\frac{\chi_{B(x,t)}(y)}{t^{n}}t^{n(1-1/r)-1}\dd t = \int_{|x-y|}^{\infty}t^{(-n/r)-1}\dd t = \frac{r}{n}|x-y|^{-n/r},
\end{equation*}
故考虑交换积分顺序,
\begin{equation}
\begin{aligned}
|I_{r}f(x)| & \leq \int_{\mathbb{R}^{n}}|x-y|^{-n/r}|f(y)|\dd y \\
& = \frac{n}{r}\int_{\mathbb{R}^{n}}\int_{0}^{\infty}\frac{\chi_{B(x,t)}(y)}{t^{n}}t^{n(1-1/r)-1}\dd t\,|f(y)|\dd y \\
& = \frac{n}{r}\int_{0}^{\infty}f_{t}(x)t^{n(1-1/r)-1}\dd t.
\end{aligned}
\end{equation}

设 $\epsilon > 0$,在完成必要的计算之后再确定其取值,将上式右端的积分分为
\begin{equation*}
\mathrm{I} := \int_{0}^{\epsilon}, \quad \mathrm{II} := \int_{\epsilon}^{\infty}
\end{equation*}
两部分分别估计.

注意 $f_{t}(x) \leq C\mathcal{M}f(x)$,
\begin{equation*}
\mathrm{I} = \int_{0}^{\epsilon}f_{t}(x)t^{n(1-1/r)-1}\dd t \leq C\mathcal{M}f(x)\epsilon^{n(1-1/r)}.
\end{equation*}

由 Hölder 不等式与 $f \in L^{p}(\mathbb{R}^{n})$,
\begin{equation}
\begin{aligned}
f_{t}(x) & = \frac{1}{t^{n}}\int_{B(x,t)}|f(y)|\dd y \\
& \leq \frac{1}{t^{n}}|B(x,t)|^{1-1/p}\cdot\left(\int_{B(x,t)}|f(y)|^{p}\dd y\right)^{1/p} \leq Ct^{-n/p}\|f\|_{p}.
\end{aligned}
\end{equation}

由 \eqref{eq:4.2.3} 推出的
\begin{equation*}
1 - \frac{1}{r} - \frac{1}{p} = -\frac{1}{q}
\end{equation*}
与上式,
\begin{equation}
\begin{aligned}
\mathrm{II} &= \int_{\epsilon}^{\infty}f_{t}(x)t^{n(1-1/r)-1}\dd t \\
&\leq C\|f\|_{p}\int_{\epsilon}^{\infty}t^{n(1-1/r-1/p)-1}\dd t \\
&= C\|f\|_{p}\epsilon^{-n/q}.
\end{aligned}
\end{equation}

因此
\begin{equation}
|I_{r}f(x)| \leq C\left(\mathcal{M}f(x)\epsilon^{n(1-1/r)} + \|f\|_{p}\epsilon^{-n/q}\right).\tag{4.2.5}
\label{eq:4.2.5}\end{equation}
此右端在满足
\begin{equation*}
\mathcal{M}f(x)\epsilon^{n(1-1/r)} = \|f\|_{p}\epsilon^{-n/q}
\end{equation*}
的 $\epsilon$ 处几乎取到最小.即,注意由 \eqref{eq:4.2.3} 推出的
\begin{equation*}
1 - \frac{1}{r} + \frac{1}{q} = \frac{1}{p}, \quad 1 - \frac{p}{q} = p - \frac{p}{r},
\end{equation*}
取
\begin{equation*}
\epsilon = \left(\frac{\|f\|_{p}}{\mathcal{M}f(x)}\right)^{p/n},
\end{equation*}
得
\begin{equation}
|I_{r}f(x)| \leq C\|f\|_{p}^{1-p/q}\mathcal{M}f(x)^{p/q}.\tag{4.2.6}
\label{eq:4.2.6}\end{equation}

由极大函数的 $L^{p}(\mathbb{R}^{n})$ 有界性,
\begin{equation*}
\int_{\mathbb{R}^{n}}(\mathcal{M}f(x))^{p}\dd x \leq C\|f\|_{p}^{p},
\end{equation*}
故
\begin{equation*}
\int_{\mathbb{R}^{n}}|I_{r}f(x)|^{q}\dd x \leq C\|f\|_{p}^{q-p}\int_{\mathbb{R}^{n}}(\mathcal{M}f(x))^{p}\dd x \leq C\|f\|_{p}^{q}.
\end{equation*}
两边开 $q$ 次方,
\begin{equation*}
\|I_{r}f\|_{q} \leq C\|f\|_{p}.
\end{equation*}
\end{proof}
